\section{\href{https://doi.org/10.1088/1402-4896/abf18b}{Double photoionization of atomic oxygen: Feshbach resonances in the two-electron continuum
}}

\subsection*{Supervisor: Prof Connor Ballance}

We describe a joint experimental and theoretical investigation on oxygen double photoionization—the
emission of two electrons from atomic oxygen following single photon absorption. High-resolution
experimental measurements were performed at the Advanced Light Source, revealing sharp resonance
structure superimposed on the more familiar Wannier-like, nearly-linear background. These resonance
features are attributed to ionization-plus-excitation Feshbach resonances embedded in the double
ionization continuum, doubly-excited states that lie above the double-ionization threshold. Such
features are absent in the double photoionization cross section of He, or other quasi-two-electron
systems, for which the doubly-ionized atomic core remains inert. For a corresponding theoretical
analysis, the R-matrix with pseudostates (RMPS) method was invoked by calculating nal-state, two-
electron resonances-plus-continua wavefunctions and corresponding single-photon absorption cross
sections. Overall agreement is found in the direct, background double photoionization cross section.
However, the RMPS method, using a small basis due to practical computational limitations, was unable
to reproduce quantitatively the smooth background or the sharper resonance features observed in the
measurements, showing instead large-scale oscillations about the experimental background, and
characteristic pseudoresonance jitter, associated with an insufficient convergence of the pseudostate
representation to the true two-electron in nite series of Feshbach resonances embedded in the two-
electron continuum. The prominent resonance structure observed highlights the need to consider
multiple excitation processes in atoms more complex than He or quasi-two-electron systems.
